\subsection{本文的安排}\label{sec:roadmap}

本文分三部。第一部處理偏誤的刻畫與診斷：第\ref{sec:method}節把標準 LC 辨識為
取平方損失的 $M$-估計量，計算其估計方程在真值處的期望，寫出該期望的一般形式，
再分別代入 Poisson、常態與過度離散三種分布；第\ref{sec:second2}節在同一個框架
下導出對數尺度的變異數、估計方程的靈敏度、校正所付的變異數代價，
以及相對於 Poisson 最大概似的效率；第\ref{sec:valid}節交代研究資料與模擬設定，
並以臺灣女性的死亡資料核對前兩節的各式；第\ref{sec:feas}節把前述各式所需的
輸入限縮為曝露數與一組粗略死亡率，據以構造一個事前的可行性檢定與分級，
並以我國 368 個鄉鎮市區的年齡結構實作。

第二部處理偏誤的校正。第\ref{sec:corr}節把第一部所得的量代回估計流程，
分解參數偏誤的確定性與隨機成分，並推導加權的最適形式；
第\ref{sec:result}節為實證結果，依序比較整體指標、逐年齡的改善剖面、
平均餘命的偏誤與離散、與參考母體修勻法的對照、模型誤設下的穩健性，
以及兩項改善得以疊加的條件。

第三部為第\ref{sec:disc}節，陳述本文修正的限制與界限，並收束全文的結論。

